\section{Executing every saved checkpoint}
\label{app:sweep}

\paragraph{Policies and hardware.}
Each run saved seven checkpoints, at updates 0, 32, \dots, 192. In every update-0 file each adapter
output matrix is exactly zero, so all six compute the frozen base model; the unweighted
seed-20260825 copy stands for them. That leaves 37 distinct policies: the 30 trained checkpoints the
rule did not select, the untrained model and the six selected checkpoints. All 37 ran on the same 61
episodes, in the same order, with the same simulator settings, step cap and tie rule as the first
closed-loop evaluation, on NVIDIA L40S GPUs (driver 610.43.02), two processes per GPU, in an Apptainer image of
NVIDIA's PyTorch 25.10 container (torch 2.9.0a0, transformers 4.57.6, peft 0.18.1, habitat-sim 0.3.3). The
preregistration was frozen in version control on 9 October 2026, before any unselected checkpoint
was executed or scored. It was not filed with a public registry, and the frozen text is in the code
release. Every result records the frozen version, the GPU's name and the container image's hash. Table~\ref{tab:sweep} gives every policy's \textsc{stop} presses and oracle successes. All
$2{,}257$ episode runs completed, and none needed a retry after an error.

%
\begin{table}[t]
\centering
\small
\caption{Every saved checkpoint executed on the 61 main episodes on NVIDIA L40S GPUs:
\textsc{stop} presses, with oracle successes in parentheses. Bold: the selection rule's pick, re-run
here. Every trained checkpoint succeeds in $0/61$ episodes: exact two-sided $95\%$ interval
$0$--$5.9\%$ for one checkpoint alone, and at most $10.3\%$ when all 37 are bounded together
(one-sided, Bonferroni). Update 0, one untrained model shared by all six runs, succeeds in
$1/61$ (bounded together: at most $13.7\%$), enters the radius in 1 and presses
\textsc{stop} 50 times.}
\label{tab:sweep}
\begin{tabular}{@{}llrrrrrr@{}}
\toprule
 & & \multicolumn{6}{c}{Update} \\
\cmidrule(l){3-8}
Treatment & Seed & 32 & 64 & 96 & 128 & 160 & 192 \\
\midrule
Unweighted & 20260825 & 0 (1) & 3 (1) & 0 (1) & 0 (1) & \textbf{0 (1)} & 0 (1) \\
 & 20260826 & 8 (0) & 5 (1) & 0 (1) & 0 (1) & \textbf{0 (1)} & 0 (1) \\
 & 20260827 & 7 (1) & 0 (1) & 0 (1) & 0 (1) & \textbf{0 (1)} & 0 (0) \\
\midrule
Weighted & 20260825 & 3 (0) & 6 (0) & 0 (1) & 4 (0) & 8 (1) & \textbf{5 (1)} \\
 & 20260826 & 22 (0) & 42 (0) & 0 (0) & \textbf{0 (1)} & 1 (0) & 6 (1) \\
 & 20260827 & 25 (0) & 0 (0) & 0 (2) & 0 (1) & 15 (0) & \textbf{3 (0)} \\
\bottomrule
\end{tabular}
\end{table}


\paragraph{Pre-run check.}
Before the preregistration was frozen, the weighted seed-20260825 checkpoint at update 192 was scored offline and
executed on eight \texttt{val\_seen} episodes in neither suite (episodes 4, 5, 10, 11, 13, 14, 16
and 17): alone, two at once on one GPU, and across a forced stop and resume. The same job replayed
the expert and an always-forward policy on all 61 main and 61 confirmation episodes. Expert replay
ended within $3.2\times10^{-6}$ m of the first platform's final distance on every main episode, the offline
scores reproduced the training records within the preregistered bars, and two processes kept the
GPU $94\%$ busy. An earlier run also passed; the freeze relies on a repeat after later code changes.

\paragraph{Offline scores.}
As the first job after the freeze, all 37 policies were scored on the 256 selection, 640 natural
and 256 balanced decisions, before any rollout was read. Of the $14{,}848$ decisions training had
also recorded, 90 are bit-identical; the largest change in a mean cross-entropy is $0.043$ nats,
within the preregistered $0.05$, and no decision won by a nat or more changed its action.

\paragraph{The first comparison: a better checkpoint.}
For each of the 31 candidates the null hypothesis is that it succeeds on no more episodes than its
comparator: its own run's selected checkpoint, and for the untrained model each of the six selected
checkpoints, all of which it must beat. With $b$ the episodes only the candidate solves and $c$ those
only its comparator solves, $p=P(\mathrm{Binomial}(b+c,1/2)\ge b)$, rejected under the
alpha-consistent form of Tarone's screen \citep{tarone1990modified,hommel1998improvements}: with
every comparator at zero successes, a candidate needs at least 5 of 61 when it is the only one that
could reach significance. Every trained candidate has $b=c=0$, so $p=1$; the untrained model has
$b=1$, $c=0$ against each selected checkpoint, so $p=0.5$. No candidate clears the bar, so
none goes on to a confirmation run.

\paragraph{The second comparison: choosing by balanced loss.}
This rule takes, within each run, the checkpoint with the lowest balanced-set cross-entropy, ties to the
earlier update. It was named before any unselected checkpoint was executed or scored, though after
the selected checkpoints' balanced scores were known. The balanced-loss rule picks updates 64, 192 and 160 in the
unweighted runs and 160, 192 and 160 in the weighted runs, the selection rule's own checkpoint in
one run of six. Its picks solve no episode, nor do the selection rule's, so $b=c=0$ and $p=1$: the
test fails, and the balanced loss is reported as a warning, not as a way to choose. Keeping each
run's last checkpoint solves none either, so it did as well as the balanced-loss rule, and class balance is not shown
to matter here. Two more rules are
reported only as description, with unadjusted $p$-values.
The lowest selection loss among checkpoints that pass the existence gate abstains in all six runs,
because none passes. The highest strict balanced \textsc{stop} recall, ties to the lower balanced
loss, picks the untrained model in the three unweighted runs and so solves one episode, 358, where
the selection rule solves none ($p=0.5$, unadjusted, descriptive); this is not a test, and we do not
call that rule better.

\paragraph{Selection regret.}
In each run, the most successes among its seven checkpoints minus the selected checkpoint's is $1$,
in all six runs, from the untrained model's episode 358. This is biased upward: it is the largest of
seven counts on the same 61 episodes.

\paragraph{The offline \textsc{stop} count.}
All 18 trained unweighted checkpoints predict \textsc{stop} on none of the 256 selection decisions.
The three selected ones pressed it in no episode on the GB10, and again in none here. We predicted
that the 15 unselected ones would press it in no episode either. The prediction failed at update 64
of seed 20260825 (3 presses), updates 32 and 64 of seed 20260826 (8 and 5) and update 32 of seed
20260827 (7): $23$ presses, all at the first step, $6.6$--$12.9$ m from the goal.

\paragraph{Offline criteria against outcomes.}
No criterion can be associated with success, which only one policy reached. Within runs, checkpoints
with a higher selection cross-entropy pressed \textsc{stop} more often (mean within-run Spearman
$0.56$, permutation $p=0.0011$, below the preregistered $0.0025$); the association with oracle
success does not reach that threshold ($p=0.0027$). The existence gate failed at all 37 policies; balanced \textsc{stop}
recall was at least 16 of 64 at three of them, none of which succeeded.

\paragraph{The two platforms.}
The six selected checkpoints ran on the GB10 first and on the L40S here. Of their 366 episodes, 347
have identical action sequences and 364 identical success, oracle-success and stop flags; each of
the other 19 first differs at a decision the GB10 won by $0.125$ nats or less, so the two machines
agree up to near-ties. The two flag differences are in one weighted checkpoint: one episode ends in
\textsc{stop} on the L40S only, and another enters the goal radius on the GB10 only; neither succeeds
on either. The comparison cannot say why two rollouts differ: the GPU, its kernels, the renderer and
the runner code all differ at once.

